% page 565 of the facsimile (image p-564 of the scan)
% block 165.1 x 233.3 mm ; left margin 36.9 mm
\pgc{15.240mm}{0.000mm}{
\l{\cel{50}557.}
\l{}
\l{\sou{strucho}. (\sou{zool.}) Stelt-ucelo, kun longa kolo, ali rudimenta, e}
\l{\cel{3}la ucelo maxim granda, remarkinda pro lua devoremeso e la}
\l{\cel{3}forteso di lua digesto-sistemo. - L.\cel{1}\sou{struthio camelus}\cel{1}-DEFI.}
\l{}
\l{\sou{strukturo}. I. Maniero ye qua edifico esas konstruktita. - II.}
\l{\cel{3}(\sou{metaf.}) Maniero ye qua esas dispozita la parti di ensemblo.-}
\l{\cel{3}DEFIRS.}
\l{}
\l{\sou{strumo}. (\sou{patol.}) Tumoro sen-dolora, ye la parto avana di la}
\l{\cel{3}kolo. - F.}
\l{}
\l{\sou{studento}. La adolecanto qua lernas en universitato, en fakultato.}
\l{\cel{3}- DEFIRS.}
\l{}
\l{\sou{studiar}. (\sou{trans. e netrans.}) I. Explorar atencoze, minucioze}
\l{\cel{3}(ulo) por determinar olua karaktero, signifiko. - II. (\sou{metaf.})}
\l{\cel{3}Explorar (la cirkonstanci, la kunjunteso di la eventi) ante}
\l{\cel{3}agar. - III. Explorar atencoze e minucioze (verko) por inter-}
\l{\cel{3}tar lu, imitar lu, e c. - DEFIS.\cel{25}pre-}
\l{}
\l{\sou{stufar}. (\sou{trans.}) Koquar (karno-bloko, pultruno) en recipiento}
\l{\cel{3}klozita (+kluza) por impedar la eskapo di la vaporo e di la}
\l{\cel{3}odoro - koquo kun acesoraji (fungi, keneli, e c.) - EFIS.}
\l{}
\l{\sou{stuko}. Induturo qua imitas marmoro, quan onu aplikas sur muro,}
\l{\cel{3}sur kolono, e quan onu polisas kande lu hardeskis per sik-}
\l{\cel{3}esko. - DEFIRS.}
\l{}
\l{\sou{stulo}. Sidilo, di qua la dors-apogilo esas min larja kam olta}
\l{\cel{3}di fotelo, e sen brakii.- DER.}
\l{}
\l{\sou{stulta}. Rido-mokinda pro judiko-defekto quan lu unike ne per-}
\l{\cel{3}ceptas, ne koncias. - I.}
\l{}
\l{\sou{stumpo}. I. (\sou{kirurg.}) La parto qua restas de membro amputita. -}
\l{\cel{3}II. La centro di ula frukti, la pedo di ula legumi de qui}
\l{\cel{3}onu deprenis la tota parto manjebla. -\cel{2}DE.}
\l{}
\l{\sou{stuntar}. (\sou{trans.})\cel{2}Haltigar (lua) kresko normala. - E.}
\l{}
\l{\sou{stupo}. La parto maxim grosiera ek la filamenti di kanabo o di}
\l{\cel{3}lino. - FIS.}
\l{}
\l{\sou{stupida}. Di qua la spirito, la mento, la psiko esas quaze tor-}
\l{\cel{3}poranta. - DEFIS.}
\l{}
\l{\sou{stuporar}. (\sou{netrans.}) Produktar sorto di torporo dum qua la penso-}
\l{\cel{3}fakultato aspektas quaze nihiligita. - EFIS.}
\l{}
\l{sturgo. (\sou{zool.}) Grosa fisho qua natas de maro aden la granda}
\l{\cel{3}fluvii, e di qua la ovi uzesas por facar kaviaro. - L. aci-}
\l{\cel{3}\sou{penser sturio}. - EFIS.}
}
